% flujo_experimentos.tex — Pipeline experimental (Figura en ICRv7)
% Uso:  \resizebox{\textwidth}{!}{% flujo_experimentos.tex — Pipeline experimental (Figura en ICRv7)
% Uso:  \resizebox{\textwidth}{!}{% flujo_experimentos.tex — Pipeline experimental (Figura en ICRv7)
% Uso:  \resizebox{\textwidth}{!}{% flujo_experimentos.tex — Pipeline experimental (Figura en ICRv7)
% Uso:  \resizebox{\textwidth}{!}{\input{figures/flujo_experimentos.tex}}

\begin{tikzpicture}[
    font=\small,
    >={Stealth[length=3mm, width=2mm]},
    % ── Estilos de nodos ────────────────────────────
    dataset/.style={
        rectangle, rounded corners=3pt, drop shadow,
        draw=orange!70!black, thick, fill=orange!8,
        text width=3.2cm, align=center, minimum height=1.2cm
    },
    proceso/.style={
        rectangle, rounded corners=3pt, drop shadow,
        draw=cyan!60!black, thick, fill=cyan!5,
        text width=3.2cm, align=center, minimum height=1.2cm
    },
    loranode/.style={
        rectangle, rounded corners=3pt, drop shadow,
        draw=blue!60!violet, thick, fill=blue!3!violet!2!white,
        text width=4cm, align=center, minimum height=1.4cm
    },
    escala/.style={
        rectangle, rounded corners=3pt, drop shadow,
        draw=yellow!50!orange, line width=1.5pt, fill=yellow!12,
        text width=3.8cm, align=center, minimum height=1.2cm
    },
    peso/.style={
        rectangle, rounded corners=3pt, drop shadow,
        draw=green!55!black, thick, fill=green!6,
        text width=1.5cm, align=center, minimum height=0.9cm
    },
    evalnode/.style={
        rectangle, rounded corners=3pt, drop shadow,
        draw=violet!65!black, thick, fill=violet!5,
        text width=3.2cm, align=center, minimum height=1.2cm
    },
    benchnode/.style={
        rectangle, rounded corners=3pt, drop shadow,
        draw=teal!65!black, thick, fill=teal!5,
        text width=4cm, align=center, minimum height=1.1cm
    },
    metS/.style={
        rectangle, rounded corners=3pt, drop shadow,
        draw=red!55!black, thick, fill=red!5,
        text width=3cm, align=center, minimum height=1.1cm
    },
    metR/.style={
        rectangle, rounded corners=3pt, drop shadow,
        draw=green!50!black, thick, fill=green!5,
        text width=2.8cm, align=center, minimum height=1.1cm
    },
    resultado/.style={
        rectangle, rounded corners=8pt, drop shadow,
        draw=magenta!55!black, line width=1.2pt, fill=magenta!5,
        text width=3.4cm, align=center, minimum height=1cm
    },
    % ── Estilos de flechas ──────────────────────────
    arrow/.style={->, thick, draw=gray!70},
    phasearrow/.style={->, line width=1pt, draw=gray!45, densely dashed},
    % ── Estilo de cajas de fase ─────────────────────
    fasebox/.style={rounded corners=10pt, thick, inner sep=12pt},
    faselabel/.style={font=\normalsize\bfseries}
]

% ═══════════════════════════════════════════════════
% FASE 1 — Preparación de los datos  (arriba-izq)
% ═══════════════════════════════════════════════════

\node[dataset] (D1) at (0, 9) {%
    \textbf{Dataset~I · STEM}\\[2pt]
    Pares contrafactuales en\\dominio técnico-científico};

\node[proceso] (T1) at (5, 9) {%
    Tokenización con\\[2pt]
    \textbf{prompting de género}};

\node[dataset] (D2) at (0, 6) {%
    \textbf{Dataset~II · Roles}\\
    \textbf{de género}\\[2pt]
    Facetas: liderazgo,\\resiliencia, empatía\ldots};

\node[proceso] (T2) at (5, 6) {%
    \textbf{Inyección directa}\\[2pt]
    de pares\\contrafactuales};

\draw[arrow] (D1) -- (T1);
\draw[arrow] (D2) -- (T2);

% ═══════════════════════════════════════════════════
% FASE 2 — Arquitectura y Entrenamiento  (centro)
% ═══════════════════════════════════════════════════

\node[loranode] (LORA) at (14, 4.5) {%
    \textbf{LoRA / QLoRA}\\[2pt]
    \scriptsize GPT-2 Base · Medium · XL\\[-1pt]
    \scriptsize Llama-3.2-1B · 3B Instruct\\[3pt]
    \small Módulos Attention + MLP\\[1pt]
    \textit{Pesos base congelados}};

\node[escala] (ESCALA) at (14, 1.8) {%
    \textbf{Escalado Paramétrico}\\[2pt]
    \scriptsize $2\;\text{corpora}\times 5\;\text{modelos}\times 3\;\text{niveles}$\\[1pt]
    \small$=\;\mathbf{30}$ \textbf{entrenamientos}};

\node[peso] (P1) at (11, -0.5) {\textbf{1\,\%}\\Pesos};
\node[peso] (P2) at (14, -0.5) {\textbf{2\,\%}\\Pesos};
\node[peso] (P3) at (17, -0.5) {\textbf{3\,\%}\\Pesos};

\draw[arrow] (LORA)  -- (ESCALA);
\draw[arrow] (ESCALA) -- node[left,  font=\scriptsize, pos=0.6] {${\approx}\!1\%$} (P1);
\draw[arrow] (ESCALA) -- (P2);
\draw[arrow] (ESCALA) -- node[right, font=\scriptsize, pos=0.6] {${\approx}\!3\%$} (P3);

% ═══════════════════════════════════════════════════
% FASE 3 — Evaluación y Contraste  (abajo-der)
% ═══════════════════════════════════════════════════

\node[evalnode] (VS) at (25, -1.5) {%
    Modelo con adaptador\\[2pt]
    \textbf{vs.}\\[2pt]
    Modelo base sin LoRA};

\node[benchnode] (BENCH) at (25, -3.8) {%
    \textbf{Benchmarks}\\[2pt]
    StereoSet · CrowS-Pairs · BOLD};

\node[metS] (SESGO) at (22.2, -6) {%
    \textbf{Métricas de Sesgo}\\[2pt]
    Stereotype Score\\PLL · Regard};

\node[metR] (REND) at (27.8, -6) {%
    \textbf{Métricas de}\\
    \textbf{Rendimiento}\\[2pt]
    LMS · ICAT};

\node[resultado] (TRADEOFF) at (25, -8.2) {%
    \textbf{Análisis del trade-off}\\[2pt]
    equidad $\Leftrightarrow$ rendimiento};

\draw[arrow] (VS)    -- (BENCH);
\draw[arrow] (BENCH)  -- (SESGO);
\draw[arrow] (BENCH)  -- (REND);
\draw[arrow] (SESGO)  -- (TRADEOFF);
\draw[arrow] (REND)   -- (TRADEOFF);

% ═══════════════════════════════════════════════════
% Flechas inter-fase  (diagonales punteadas)
% ═══════════════════════════════════════════════════

\draw[phasearrow] (T1.east)  to[out=-10, in=150] (LORA.north west);
\draw[phasearrow] (T2.east)  to[out=-5,  in=170] (LORA.west);

\draw[phasearrow] (P1.south) to[out=-30, in=160] (VS.north west);
\draw[phasearrow] (P2.south) to[out=-25, in=140] (VS.north);
\draw[phasearrow] (P3.east)  to[out=-10, in=115] (VS.north east);

% ═══════════════════════════════════════════════════
% Cajas de fase  (capa de fondo)
% ═══════════════════════════════════════════════════

\begin{pgfonlayer}{background}
    \node[fasebox, fit=(D1)(T1)(D2)(T2),
          fill=blue!4, draw=blue!45!black] (F1box) {};

    \node[fasebox, fit=(LORA)(ESCALA)(P1)(P2)(P3),
          fill=green!4, draw=green!45!black] (F2box) {};

    \node[fasebox, fit=(VS)(BENCH)(SESGO)(REND)(TRADEOFF),
          fill=violet!4, draw=violet!40!black] (F3box) {};
\end{pgfonlayer}

% ═══════════════════════════════════════════════════
% Etiquetas de fase
% ═══════════════════════════════════════════════════

\node[faselabel, blue!55!black, above=3pt of F1box.north]
    {\textsc{Fase~1} --- Preparación de los datos};

\node[faselabel, green!45!black, above=3pt of F2box.north]
    {\textsc{Fase~2} --- Arquitectura y Entrenamiento};

\node[faselabel, violet!50!black, above=3pt of F3box.north]
    {\textsc{Fase~3} --- Evaluación y Contraste};

\end{tikzpicture}
}

\begin{tikzpicture}[
    font=\small,
    >={Stealth[length=3mm, width=2mm]},
    % ── Estilos de nodos ────────────────────────────
    dataset/.style={
        rectangle, rounded corners=3pt, drop shadow,
        draw=orange!70!black, thick, fill=orange!8,
        text width=3.2cm, align=center, minimum height=1.2cm
    },
    proceso/.style={
        rectangle, rounded corners=3pt, drop shadow,
        draw=cyan!60!black, thick, fill=cyan!5,
        text width=3.2cm, align=center, minimum height=1.2cm
    },
    loranode/.style={
        rectangle, rounded corners=3pt, drop shadow,
        draw=blue!60!violet, thick, fill=blue!3!violet!2!white,
        text width=4cm, align=center, minimum height=1.4cm
    },
    escala/.style={
        rectangle, rounded corners=3pt, drop shadow,
        draw=yellow!50!orange, line width=1.5pt, fill=yellow!12,
        text width=3.8cm, align=center, minimum height=1.2cm
    },
    peso/.style={
        rectangle, rounded corners=3pt, drop shadow,
        draw=green!55!black, thick, fill=green!6,
        text width=1.5cm, align=center, minimum height=0.9cm
    },
    evalnode/.style={
        rectangle, rounded corners=3pt, drop shadow,
        draw=violet!65!black, thick, fill=violet!5,
        text width=3.2cm, align=center, minimum height=1.2cm
    },
    benchnode/.style={
        rectangle, rounded corners=3pt, drop shadow,
        draw=teal!65!black, thick, fill=teal!5,
        text width=4cm, align=center, minimum height=1.1cm
    },
    metS/.style={
        rectangle, rounded corners=3pt, drop shadow,
        draw=red!55!black, thick, fill=red!5,
        text width=3cm, align=center, minimum height=1.1cm
    },
    metR/.style={
        rectangle, rounded corners=3pt, drop shadow,
        draw=green!50!black, thick, fill=green!5,
        text width=2.8cm, align=center, minimum height=1.1cm
    },
    resultado/.style={
        rectangle, rounded corners=8pt, drop shadow,
        draw=magenta!55!black, line width=1.2pt, fill=magenta!5,
        text width=3.4cm, align=center, minimum height=1cm
    },
    % ── Estilos de flechas ──────────────────────────
    arrow/.style={->, thick, draw=gray!70},
    phasearrow/.style={->, line width=1pt, draw=gray!45, densely dashed},
    % ── Estilo de cajas de fase ─────────────────────
    fasebox/.style={rounded corners=10pt, thick, inner sep=12pt},
    faselabel/.style={font=\normalsize\bfseries}
]

% ═══════════════════════════════════════════════════
% FASE 1 — Preparación de los datos  (arriba-izq)
% ═══════════════════════════════════════════════════

\node[dataset] (D1) at (0, 9) {%
    \textbf{Dataset~I · STEM}\\[2pt]
    Pares contrafactuales en\\dominio técnico-científico};

\node[proceso] (T1) at (5, 9) {%
    Tokenización con\\[2pt]
    \textbf{prompting de género}};

\node[dataset] (D2) at (0, 6) {%
    \textbf{Dataset~II · Roles}\\
    \textbf{de género}\\[2pt]
    Facetas: liderazgo,\\resiliencia, empatía\ldots};

\node[proceso] (T2) at (5, 6) {%
    \textbf{Inyección directa}\\[2pt]
    de pares\\contrafactuales};

\draw[arrow] (D1) -- (T1);
\draw[arrow] (D2) -- (T2);

% ═══════════════════════════════════════════════════
% FASE 2 — Arquitectura y Entrenamiento  (centro)
% ═══════════════════════════════════════════════════

\node[loranode] (LORA) at (14, 4.5) {%
    \textbf{LoRA / QLoRA}\\[2pt]
    \scriptsize GPT-2 Base · Medium · XL\\[-1pt]
    \scriptsize Llama-3.2-1B · 3B Instruct\\[3pt]
    \small Módulos Attention + MLP\\[1pt]
    \textit{Pesos base congelados}};

\node[escala] (ESCALA) at (14, 1.8) {%
    \textbf{Escalado Paramétrico}\\[2pt]
    \scriptsize $2\;\text{corpora}\times 5\;\text{modelos}\times 3\;\text{niveles}$\\[1pt]
    \small$=\;\mathbf{30}$ \textbf{entrenamientos}};

\node[peso] (P1) at (11, -0.5) {\textbf{1\,\%}\\Pesos};
\node[peso] (P2) at (14, -0.5) {\textbf{2\,\%}\\Pesos};
\node[peso] (P3) at (17, -0.5) {\textbf{3\,\%}\\Pesos};

\draw[arrow] (LORA)  -- (ESCALA);
\draw[arrow] (ESCALA) -- node[left,  font=\scriptsize, pos=0.6] {${\approx}\!1\%$} (P1);
\draw[arrow] (ESCALA) -- (P2);
\draw[arrow] (ESCALA) -- node[right, font=\scriptsize, pos=0.6] {${\approx}\!3\%$} (P3);

% ═══════════════════════════════════════════════════
% FASE 3 — Evaluación y Contraste  (abajo-der)
% ═══════════════════════════════════════════════════

\node[evalnode] (VS) at (25, -1.5) {%
    Modelo con adaptador\\[2pt]
    \textbf{vs.}\\[2pt]
    Modelo base sin LoRA};

\node[benchnode] (BENCH) at (25, -3.8) {%
    \textbf{Benchmarks}\\[2pt]
    StereoSet · CrowS-Pairs · BOLD};

\node[metS] (SESGO) at (22.2, -6) {%
    \textbf{Métricas de Sesgo}\\[2pt]
    Stereotype Score\\PLL · Regard};

\node[metR] (REND) at (27.8, -6) {%
    \textbf{Métricas de}\\
    \textbf{Rendimiento}\\[2pt]
    LMS · ICAT};

\node[resultado] (TRADEOFF) at (25, -8.2) {%
    \textbf{Análisis del trade-off}\\[2pt]
    equidad $\Leftrightarrow$ rendimiento};

\draw[arrow] (VS)    -- (BENCH);
\draw[arrow] (BENCH)  -- (SESGO);
\draw[arrow] (BENCH)  -- (REND);
\draw[arrow] (SESGO)  -- (TRADEOFF);
\draw[arrow] (REND)   -- (TRADEOFF);

% ═══════════════════════════════════════════════════
% Flechas inter-fase  (diagonales punteadas)
% ═══════════════════════════════════════════════════

\draw[phasearrow] (T1.east)  to[out=-10, in=150] (LORA.north west);
\draw[phasearrow] (T2.east)  to[out=-5,  in=170] (LORA.west);

\draw[phasearrow] (P1.south) to[out=-30, in=160] (VS.north west);
\draw[phasearrow] (P2.south) to[out=-25, in=140] (VS.north);
\draw[phasearrow] (P3.east)  to[out=-10, in=115] (VS.north east);

% ═══════════════════════════════════════════════════
% Cajas de fase  (capa de fondo)
% ═══════════════════════════════════════════════════

\begin{pgfonlayer}{background}
    \node[fasebox, fit=(D1)(T1)(D2)(T2),
          fill=blue!4, draw=blue!45!black] (F1box) {};

    \node[fasebox, fit=(LORA)(ESCALA)(P1)(P2)(P3),
          fill=green!4, draw=green!45!black] (F2box) {};

    \node[fasebox, fit=(VS)(BENCH)(SESGO)(REND)(TRADEOFF),
          fill=violet!4, draw=violet!40!black] (F3box) {};
\end{pgfonlayer}

% ═══════════════════════════════════════════════════
% Etiquetas de fase
% ═══════════════════════════════════════════════════

\node[faselabel, blue!55!black, above=3pt of F1box.north]
    {\textsc{Fase~1} --- Preparación de los datos};

\node[faselabel, green!45!black, above=3pt of F2box.north]
    {\textsc{Fase~2} --- Arquitectura y Entrenamiento};

\node[faselabel, violet!50!black, above=3pt of F3box.north]
    {\textsc{Fase~3} --- Evaluación y Contraste};

\end{tikzpicture}
}

\begin{tikzpicture}[
    font=\small,
    >={Stealth[length=3mm, width=2mm]},
    % ── Estilos de nodos ────────────────────────────
    dataset/.style={
        rectangle, rounded corners=3pt, drop shadow,
        draw=orange!70!black, thick, fill=orange!8,
        text width=3.2cm, align=center, minimum height=1.2cm
    },
    proceso/.style={
        rectangle, rounded corners=3pt, drop shadow,
        draw=cyan!60!black, thick, fill=cyan!5,
        text width=3.2cm, align=center, minimum height=1.2cm
    },
    loranode/.style={
        rectangle, rounded corners=3pt, drop shadow,
        draw=blue!60!violet, thick, fill=blue!3!violet!2!white,
        text width=4cm, align=center, minimum height=1.4cm
    },
    escala/.style={
        rectangle, rounded corners=3pt, drop shadow,
        draw=yellow!50!orange, line width=1.5pt, fill=yellow!12,
        text width=3.8cm, align=center, minimum height=1.2cm
    },
    peso/.style={
        rectangle, rounded corners=3pt, drop shadow,
        draw=green!55!black, thick, fill=green!6,
        text width=1.5cm, align=center, minimum height=0.9cm
    },
    evalnode/.style={
        rectangle, rounded corners=3pt, drop shadow,
        draw=violet!65!black, thick, fill=violet!5,
        text width=3.2cm, align=center, minimum height=1.2cm
    },
    benchnode/.style={
        rectangle, rounded corners=3pt, drop shadow,
        draw=teal!65!black, thick, fill=teal!5,
        text width=4cm, align=center, minimum height=1.1cm
    },
    metS/.style={
        rectangle, rounded corners=3pt, drop shadow,
        draw=red!55!black, thick, fill=red!5,
        text width=3cm, align=center, minimum height=1.1cm
    },
    metR/.style={
        rectangle, rounded corners=3pt, drop shadow,
        draw=green!50!black, thick, fill=green!5,
        text width=2.8cm, align=center, minimum height=1.1cm
    },
    resultado/.style={
        rectangle, rounded corners=8pt, drop shadow,
        draw=magenta!55!black, line width=1.2pt, fill=magenta!5,
        text width=3.4cm, align=center, minimum height=1cm
    },
    % ── Estilos de flechas ──────────────────────────
    arrow/.style={->, thick, draw=gray!70},
    phasearrow/.style={->, line width=1pt, draw=gray!45, densely dashed},
    % ── Estilo de cajas de fase ─────────────────────
    fasebox/.style={rounded corners=10pt, thick, inner sep=12pt},
    faselabel/.style={font=\normalsize\bfseries}
]

% ═══════════════════════════════════════════════════
% FASE 1 — Preparación de los datos  (arriba-izq)
% ═══════════════════════════════════════════════════

\node[dataset] (D1) at (0, 9) {%
    \textbf{Dataset~I · STEM}\\[2pt]
    Pares contrafactuales en\\dominio técnico-científico};

\node[proceso] (T1) at (5, 9) {%
    Tokenización con\\[2pt]
    \textbf{prompting de género}};

\node[dataset] (D2) at (0, 6) {%
    \textbf{Dataset~II · Roles}\\
    \textbf{de género}\\[2pt]
    Facetas: liderazgo,\\resiliencia, empatía\ldots};

\node[proceso] (T2) at (5, 6) {%
    \textbf{Inyección directa}\\[2pt]
    de pares\\contrafactuales};

\draw[arrow] (D1) -- (T1);
\draw[arrow] (D2) -- (T2);

% ═══════════════════════════════════════════════════
% FASE 2 — Arquitectura y Entrenamiento  (centro)
% ═══════════════════════════════════════════════════

\node[loranode] (LORA) at (14, 4.5) {%
    \textbf{LoRA / QLoRA}\\[2pt]
    \scriptsize GPT-2 Base · Medium · XL\\[-1pt]
    \scriptsize Llama-3.2-1B · 3B Instruct\\[3pt]
    \small Módulos Attention + MLP\\[1pt]
    \textit{Pesos base congelados}};

\node[escala] (ESCALA) at (14, 1.8) {%
    \textbf{Escalado Paramétrico}\\[2pt]
    \scriptsize $2\;\text{corpora}\times 5\;\text{modelos}\times 3\;\text{niveles}$\\[1pt]
    \small$=\;\mathbf{30}$ \textbf{entrenamientos}};

\node[peso] (P1) at (11, -0.5) {\textbf{1\,\%}\\Pesos};
\node[peso] (P2) at (14, -0.5) {\textbf{2\,\%}\\Pesos};
\node[peso] (P3) at (17, -0.5) {\textbf{3\,\%}\\Pesos};

\draw[arrow] (LORA)  -- (ESCALA);
\draw[arrow] (ESCALA) -- node[left,  font=\scriptsize, pos=0.6] {${\approx}\!1\%$} (P1);
\draw[arrow] (ESCALA) -- (P2);
\draw[arrow] (ESCALA) -- node[right, font=\scriptsize, pos=0.6] {${\approx}\!3\%$} (P3);

% ═══════════════════════════════════════════════════
% FASE 3 — Evaluación y Contraste  (abajo-der)
% ═══════════════════════════════════════════════════

\node[evalnode] (VS) at (25, -1.5) {%
    Modelo con adaptador\\[2pt]
    \textbf{vs.}\\[2pt]
    Modelo base sin LoRA};

\node[benchnode] (BENCH) at (25, -3.8) {%
    \textbf{Benchmarks}\\[2pt]
    StereoSet · CrowS-Pairs · BOLD};

\node[metS] (SESGO) at (22.2, -6) {%
    \textbf{Métricas de Sesgo}\\[2pt]
    Stereotype Score\\PLL · Regard};

\node[metR] (REND) at (27.8, -6) {%
    \textbf{Métricas de}\\
    \textbf{Rendimiento}\\[2pt]
    LMS · ICAT};

\node[resultado] (TRADEOFF) at (25, -8.2) {%
    \textbf{Análisis del trade-off}\\[2pt]
    equidad $\Leftrightarrow$ rendimiento};

\draw[arrow] (VS)    -- (BENCH);
\draw[arrow] (BENCH)  -- (SESGO);
\draw[arrow] (BENCH)  -- (REND);
\draw[arrow] (SESGO)  -- (TRADEOFF);
\draw[arrow] (REND)   -- (TRADEOFF);

% ═══════════════════════════════════════════════════
% Flechas inter-fase  (diagonales punteadas)
% ═══════════════════════════════════════════════════

\draw[phasearrow] (T1.east)  to[out=-10, in=150] (LORA.north west);
\draw[phasearrow] (T2.east)  to[out=-5,  in=170] (LORA.west);

\draw[phasearrow] (P1.south) to[out=-30, in=160] (VS.north west);
\draw[phasearrow] (P2.south) to[out=-25, in=140] (VS.north);
\draw[phasearrow] (P3.east)  to[out=-10, in=115] (VS.north east);

% ═══════════════════════════════════════════════════
% Cajas de fase  (capa de fondo)
% ═══════════════════════════════════════════════════

\begin{pgfonlayer}{background}
    \node[fasebox, fit=(D1)(T1)(D2)(T2),
          fill=blue!4, draw=blue!45!black] (F1box) {};

    \node[fasebox, fit=(LORA)(ESCALA)(P1)(P2)(P3),
          fill=green!4, draw=green!45!black] (F2box) {};

    \node[fasebox, fit=(VS)(BENCH)(SESGO)(REND)(TRADEOFF),
          fill=violet!4, draw=violet!40!black] (F3box) {};
\end{pgfonlayer}

% ═══════════════════════════════════════════════════
% Etiquetas de fase
% ═══════════════════════════════════════════════════

\node[faselabel, blue!55!black, above=3pt of F1box.north]
    {\textsc{Fase~1} --- Preparación de los datos};

\node[faselabel, green!45!black, above=3pt of F2box.north]
    {\textsc{Fase~2} --- Arquitectura y Entrenamiento};

\node[faselabel, violet!50!black, above=3pt of F3box.north]
    {\textsc{Fase~3} --- Evaluación y Contraste};

\end{tikzpicture}
}

\begin{tikzpicture}[
    font=\small,
    >={Stealth[length=3mm, width=2mm]},
    % ── Estilos de nodos ────────────────────────────
    dataset/.style={
        rectangle, rounded corners=3pt, drop shadow,
        draw=orange!70!black, thick, fill=orange!8,
        text width=3.2cm, align=center, minimum height=1.2cm
    },
    proceso/.style={
        rectangle, rounded corners=3pt, drop shadow,
        draw=cyan!60!black, thick, fill=cyan!5,
        text width=3.2cm, align=center, minimum height=1.2cm
    },
    loranode/.style={
        rectangle, rounded corners=3pt, drop shadow,
        draw=blue!60!violet, thick, fill=blue!3!violet!2!white,
        text width=4cm, align=center, minimum height=1.4cm
    },
    escala/.style={
        rectangle, rounded corners=3pt, drop shadow,
        draw=yellow!50!orange, line width=1.5pt, fill=yellow!12,
        text width=3.8cm, align=center, minimum height=1.2cm
    },
    peso/.style={
        rectangle, rounded corners=3pt, drop shadow,
        draw=green!55!black, thick, fill=green!6,
        text width=1.5cm, align=center, minimum height=0.9cm
    },
    evalnode/.style={
        rectangle, rounded corners=3pt, drop shadow,
        draw=violet!65!black, thick, fill=violet!5,
        text width=3.2cm, align=center, minimum height=1.2cm
    },
    benchnode/.style={
        rectangle, rounded corners=3pt, drop shadow,
        draw=teal!65!black, thick, fill=teal!5,
        text width=4cm, align=center, minimum height=1.1cm
    },
    metS/.style={
        rectangle, rounded corners=3pt, drop shadow,
        draw=red!55!black, thick, fill=red!5,
        text width=3cm, align=center, minimum height=1.1cm
    },
    metR/.style={
        rectangle, rounded corners=3pt, drop shadow,
        draw=green!50!black, thick, fill=green!5,
        text width=2.8cm, align=center, minimum height=1.1cm
    },
    resultado/.style={
        rectangle, rounded corners=8pt, drop shadow,
        draw=magenta!55!black, line width=1.2pt, fill=magenta!5,
        text width=3.4cm, align=center, minimum height=1cm
    },
    % ── Estilos de flechas ──────────────────────────
    arrow/.style={->, thick, draw=gray!70},
    phasearrow/.style={->, line width=1pt, draw=gray!45, densely dashed},
    % ── Estilo de cajas de fase ─────────────────────
    fasebox/.style={rounded corners=10pt, thick, inner sep=12pt},
    faselabel/.style={font=\normalsize\bfseries}
]

% ═══════════════════════════════════════════════════
% FASE 1 — Preparación de los datos  (arriba-izq)
% ═══════════════════════════════════════════════════

\node[dataset] (D1) at (0, 9) {%
    \textbf{Dataset~I · STEM}\\[2pt]
    Pares contrafactuales en\\dominio técnico-científico};

\node[proceso] (T1) at (5, 9) {%
    Tokenización con\\[2pt]
    \textbf{prompting de género}};

\node[dataset] (D2) at (0, 6) {%
    \textbf{Dataset~II · Roles}\\
    \textbf{de género}\\[2pt]
    Facetas: liderazgo,\\resiliencia, empatía\ldots};

\node[proceso] (T2) at (5, 6) {%
    \textbf{Inyección directa}\\[2pt]
    de pares\\contrafactuales};

\draw[arrow] (D1) -- (T1);
\draw[arrow] (D2) -- (T2);

% ═══════════════════════════════════════════════════
% FASE 2 — Arquitectura y Entrenamiento  (centro)
% ═══════════════════════════════════════════════════

\node[loranode] (LORA) at (14, 4.5) {%
    \textbf{LoRA / QLoRA}\\[2pt]
    \scriptsize GPT-2 Base · Medium · XL\\[-1pt]
    \scriptsize Llama-3.2-1B · 3B Instruct\\[3pt]
    \small Módulos Attention + MLP\\[1pt]
    \textit{Pesos base congelados}};

\node[escala] (ESCALA) at (14, 1.8) {%
    \textbf{Escalado Paramétrico}\\[2pt]
    \scriptsize $2\;\text{corpora}\times 5\;\text{modelos}\times 3\;\text{niveles}$\\[1pt]
    \small$=\;\mathbf{30}$ \textbf{entrenamientos}};

\node[peso] (P1) at (11, -0.5) {\textbf{1\,\%}\\Pesos};
\node[peso] (P2) at (14, -0.5) {\textbf{2\,\%}\\Pesos};
\node[peso] (P3) at (17, -0.5) {\textbf{3\,\%}\\Pesos};

\draw[arrow] (LORA)  -- (ESCALA);
\draw[arrow] (ESCALA) -- node[left,  font=\scriptsize, pos=0.6] {${\approx}\!1\%$} (P1);
\draw[arrow] (ESCALA) -- (P2);
\draw[arrow] (ESCALA) -- node[right, font=\scriptsize, pos=0.6] {${\approx}\!3\%$} (P3);

% ═══════════════════════════════════════════════════
% FASE 3 — Evaluación y Contraste  (abajo-der)
% ═══════════════════════════════════════════════════

\node[evalnode] (VS) at (25, -1.5) {%
    Modelo con adaptador\\[2pt]
    \textbf{vs.}\\[2pt]
    Modelo base sin LoRA};

\node[benchnode] (BENCH) at (25, -3.8) {%
    \textbf{Benchmarks}\\[2pt]
    StereoSet · CrowS-Pairs · BOLD};

\node[metS] (SESGO) at (22.2, -6) {%
    \textbf{Métricas de Sesgo}\\[2pt]
    Stereotype Score\\PLL · Regard};

\node[metR] (REND) at (27.8, -6) {%
    \textbf{Métricas de}\\
    \textbf{Rendimiento}\\[2pt]
    LMS · ICAT};

\node[resultado] (TRADEOFF) at (25, -8.2) {%
    \textbf{Análisis del trade-off}\\[2pt]
    equidad $\Leftrightarrow$ rendimiento};

\draw[arrow] (VS)    -- (BENCH);
\draw[arrow] (BENCH)  -- (SESGO);
\draw[arrow] (BENCH)  -- (REND);
\draw[arrow] (SESGO)  -- (TRADEOFF);
\draw[arrow] (REND)   -- (TRADEOFF);

% ═══════════════════════════════════════════════════
% Flechas inter-fase  (diagonales punteadas)
% ═══════════════════════════════════════════════════

\draw[phasearrow] (T1.east)  to[out=-10, in=150] (LORA.north west);
\draw[phasearrow] (T2.east)  to[out=-5,  in=170] (LORA.west);

\draw[phasearrow] (P1.south) to[out=-30, in=160] (VS.north west);
\draw[phasearrow] (P2.south) to[out=-25, in=140] (VS.north);
\draw[phasearrow] (P3.east)  to[out=-10, in=115] (VS.north east);

% ═══════════════════════════════════════════════════
% Cajas de fase  (capa de fondo)
% ═══════════════════════════════════════════════════

\begin{pgfonlayer}{background}
    \node[fasebox, fit=(D1)(T1)(D2)(T2),
          fill=blue!4, draw=blue!45!black] (F1box) {};

    \node[fasebox, fit=(LORA)(ESCALA)(P1)(P2)(P3),
          fill=green!4, draw=green!45!black] (F2box) {};

    \node[fasebox, fit=(VS)(BENCH)(SESGO)(REND)(TRADEOFF),
          fill=violet!4, draw=violet!40!black] (F3box) {};
\end{pgfonlayer}

% ═══════════════════════════════════════════════════
% Etiquetas de fase
% ═══════════════════════════════════════════════════

\node[faselabel, blue!55!black, above=3pt of F1box.north]
    {\textsc{Fase~1} --- Preparación de los datos};

\node[faselabel, green!45!black, above=3pt of F2box.north]
    {\textsc{Fase~2} --- Arquitectura y Entrenamiento};

\node[faselabel, violet!50!black, above=3pt of F3box.north]
    {\textsc{Fase~3} --- Evaluación y Contraste};

\end{tikzpicture}
